\documentclass[10pt,twocolumn,a4paper]{article}

% ============================================================================
% PACKAGES STANDARDS (COMPATIBILITE UNIVERSELLE TEXLIVE 2026)
% ============================================================================
\usepackage[utf8]{inputenc}
\usepackage[T1]{fontenc}
\usepackage{lmodern}
\usepackage[french,english]{babel}
\usepackage{amsmath,amssymb,amsfonts,amsthm,mathtools}
\usepackage{bm}
\usepackage{geometry}
\usepackage{booktabs}
\usepackage{array}
\usepackage{microtype}
\usepackage{xcolor}
\usepackage{graphicx}
\usepackage{listings}
\usepackage{hyperref}
\usepackage{cite}
\usepackage{titlesec}
\usepackage{cuted} % Permet l'environnement strip (équivalent widetext de revtex4)
\usepackage{tabularx}

% Configuration de la géométrie (Style Physical Review D / IPI Letters)
\geometry{
    a4paper,
    top=1.8cm,
    bottom=2.0cm,
    left=1.5cm,
    right=1.5cm,
    columnsep=0.65cm
}

% Tolérance d'espacement pour éviter tout Overfull \hbox en mode 2 colonnes
\emergencystretch=1.5em
\allowdisplaybreaks

\hypersetup{
    colorlinks=true,
    linkcolor=blue!70!black,
    citecolor=blue!70!black,
    urlcolor=blue!70!black,
    pdftitle={Holographie Celeste, Temps Thermique et Cosmologie Quantique du Consensus de Nakamoto},
    pdfauthor={Pascal Ranaora}
}

% Configuration des blocs de code (Lean 4 et Python) avec support UTF-8 mathématique
\lstset{
    basicstyle=\ttfamily\scriptsize,
    breaklines=true,
    frame=single,
    numbers=left,
    numberstyle=\tiny\color{gray},
    keywordstyle=\color{blue!75!black}\bfseries,
    commentstyle=\color{green!50!black}\itshape,
    stringstyle=\color{red!60!black},
    showstringspaces=false,
    tabsize=2,
    literate=
      {ℝ}{{$\mathbb{R}$}}1
      {ℂ}{{$\mathbb{C}$}}1
      {ℕ}{{$\mathbb{N}$}}1
      {∀}{{$\forall$}}1
      {∃}{{$\exists$}}1
      {→}{{$\to$}}1
      {←}{{$\leftarrow$}}1
      {↔}{{$\leftrightarrow$}}1
      {∑}{{$\sum$}}1
      {ω}{{$\omega$}}1
      {σ}{{$\sigma$}}1
      {ρ}{{$\rho$}}1
      {β}{{$\beta$}}1
      {γ}{{$\gamma$}}1
      {Γ}{{$\Gamma$}}1
      {η}{{$\eta$}}1
      {∞}{{$\infty$}}1
      {₁}{{$_1$}}1
      {₂}{{$_2$}}1
      {⁰}{{$^0$}}1
      {¹}{{$^1$}}1
      {²}{{$^2$}}1
      {⁻}{{$^-$}}1
      {⁺}{{$^+$}}1
      {é}{{\'e}}1
      {è}{{\`e}}1
      {ê}{{\^e}}1
      {à}{{\`a}}1
      {ù}{{\`u}}1
      {î}{{\^i}}1
      {ô}{{\^o}}1
      {ç}{{\c{c}}}1
      {É}{{\'E}}1
      {À}{{\`A}}1
}

% Formatage compact des sections
\titleformat{\section}{\large\bfseries\uppercase}{\thesection.}{0.5em}{}
\titleformat{\subsection}{\normalsize\bfseries}{\thesubsection.}{0.5em}{}
\titleformat{\subsubsection}{\normalsize\itshape}{\thesubsubsection.}{0.5em}{}

% Environnement widetext (style Physical Review D) pour les longues équations
\newenvironment{widetext}{%
  \begin{strip}
  \rule{\dimexpr(0.5\textwidth-0.5\columnsep-0.4pt)}{0.4pt}%
  \rule{0.4pt}{6pt}\par\vspace{-2pt}%
}{%
  \par\vspace{-2pt}\hfill%
  \rule[-6pt]{0.4pt}{6.4pt}%
  \rule{\dimexpr(0.5\textwidth-0.5\columnsep-0.4pt)}{0.4pt}%
  \end{strip}
}

% ============================================================================
% ENVIRONNEMENTS MATHEMATIQUES (THEOREMES, AXIOMES, DEMONSTRATIONS)
% ============================================================================
\theoremstyle{plain}
\newtheorem{theorem}{Théorème}[section]
\newtheorem{lemma}[theorem]{Lemme}
\newtheorem{proposition}[theorem]{Proposition}
\newtheorem{corollary}[theorem]{Corollaire}

\theoremstyle{definition}
\newtheorem{definition}[theorem]{Définition}
\newtheorem{axiom}{Axiome Fondamental}

\theoremstyle{remark}
\newtheorem{remark}[theorem]{Remarque}

% Commandes personnalisées pour la physique mathématique et Nakamoto
\newcommand{\kNak}{k_{\text{Nak}}}
\newcommand{\GNak}{G_{\text{Nak}}}
\newcommand{\lNak}{\ell_{\text{Nak}}}
\newcommand{\LambdaNak}{\Lambda_{\text{Nak}}}
\newcommand{\scri}{\mathcal{I}}
\newcommand{\Tr}{\operatorname{Tr}}
\newcommand{\Diff}{\operatorname{Diff}}
\newcommand{\Inv}{\operatorname{Inv}}
\newcommand{\Aut}{\operatorname{Aut}}
\newcommand{\uRet}{u_{\text{ret}}}
\newcommand{\aScale}{a_{\text{scale}}}

\begin{document}
\selectlanguage{french}

% ============================================================================
% TITRE ET RESUME PLEINE LARGEUR
% ============================================================================
\twocolumn[%
\begin{minipage}{\textwidth}
    \centering
    \vspace{-0.4cm}
    {\Large\bfseries Holographie Céleste, Temps Thermique Modulaire et Cosmologie Quantique :\\[0.2em]
    Un Dictionnaire Conditionnel entre le Consensus de Nakamoto et la Physique de l'Information sous l'Hypothèse de Simulation\par}
    \vspace{0.4cm}
    {\normalsize \textbf{Pascal Ranaora}$^{1,*}$ \quad \par}
    \vspace{0.15cm}
    {\small $^{1}$\textit{Information Physics Institute (IPI), Sydney, NSW, Australia}\\
    $^{*}$\texttt{pascal.ranaora@informationphysicsinstitute.net}\par}
    \vspace{0.15cm}
    {\small \today\par}
    \vspace{0.35cm}
    \begin{minipage}{0.92\textwidth}
        \small
        \noindent \textbf{Résumé —} Nous formulons une théorie unifiée de la gravité quantique, de l'holographie céleste et de la thermodynamique statistique hors d'équilibre, projetée par isomorphisme exact sur l'architecture du consensus distribué de Satoshi Nakamoto. En posant l'axiome fondamental d'équivalence entre la constante de Nakamoto $\kNak$ (quantum irréductible d'action cryptographique SHA-256) et la constante de Planck $h$, nous proposons de modéliser la \textit{Timechain} comme une variété d'espace-temps discrète émergeant d'un bain thermique dissipatif. Premièrement, sur la frontière asymptotique nulle ($\scri^+$) du réseau pair-à-pair (P2P), nous proposons de décrire l'invariance du registre sous les délais de propagation asynchrones par analogie avec le groupe de Bondi-Metzner-Sachs ($\text{BMS}_4$). Le théorème des gravitons mous de Weinberg appliqué au protocole \textit{gossip} de la Mempool implique l'identité de Ward de conservation monétaire stricte, tandis que la validation d'un bloc induit un effet de mémoire gravitationnelle non-linéaire de Christodoulou qui encode l'immuabilité de l'ensemble UTXO. Deuxièmement, à partir de l'algèbre de von Neumann non-commutative des observables du registre, nous dérivons l'opérateur modulaire de Tomita-Takesaki $\Delta = \rho \otimes \rho^{-1}$ et montrons formellement, en supposant l'existence d'un état fidèle, que le flot du temps thermique complexe $z = a + ib$ satisfait la condition d'équilibre de Kubo-Martin-Schwinger (KMS), où l'ajustement homéostatique de la difficulté agit comme un groupe de renormalisation ramenant en moyenne le pas temporel vers $\tau_0 = 600\,\text{s}$ (point fixe statistique ; les temps de bloc individuels restent distribués exponentiellement). Troisièmement, nous modélisons la topologie microscopique du graphe UTXO comme un réseau de spin $SU(2)$ dont l'évolution en mousse de spin (\textit{Spinfoam}) est soumise à l'équation maîtresse quantique ouverte de Gorini-Kossakowski-Sudarshan-Lindblad (GKSL). Nous montrons qu'une modélisation de type GKSL, dont la restriction diagonale se réduit à une marche aléatoire classique, retrouve la loi de ruine du joueur du livre blanc de 2008 comme état stationnaire asymptotique de Lindblad. Enfin, nous établissons la cosmologie émergente de la Timechain : l'addition séquentielle de blocs est régie par la loi d'expansion holographique de Padmanabhan $d\mathcal{V}/ds \propto (N_{\text{sur}} - N_{\text{bulk}})$, la conjecture \textit{Complexité = Volume} de Susskind, l'hypothèse de courbure de Weyl nulle de Penrose au Bloc Genesis, et un potentiel d'inflaton en escalier décrivant les \textit{Halvings}. L'équation d'état thermodynamique de Jacobson $\delta Q = T \delta S$ est enfin discutée comme analogie heuristique pour une mesure de Confiance Sociétale ($C_s$), et l'ensemble des théorèmes est accompagné d'une vérification formelle sous Lean~4 (noyau natif zéro-dépendance) et SymPy ainsi que des tests statistiques empiriques T1--T3.\par
        \vspace{0.2cm}
        \noindent \textbf{Mots-clés :} Holographie Céleste, Groupe BMS, Hypothèse du Temps Thermique, Algèbre de Tomita-Takesaki, Condition KMS, Mousses de Spin, Équation Maîtresse GKSL, Cosmologie Émergente, Consensus de Nakamoto, Vérification Formelle Lean 4.
    \end{minipage}
    \vspace{0.6cm}
\end{minipage}%
]

% ============================================================================
% SECTION 1 : Introduction et Axiomatique Conditionnelle
% ============================================================================
\section{Introduction et Axiomatique Conditionnelle}

Le problème de la coordination d'horloges dans un système distribué asynchrone en présence d'acteurs byzantins (malveillants) et le problème de l'émergence d'un espace-temps classique continu à partir de degrés de liberté quantiques sans horloge globale partagent la même obstruction mathématique fondamentale : l'invariance par difféomorphisme temporel. En relativité générale canonique, cette obstruction se manifeste par la contrainte hamiltonienne de Wheeler-DeWitt :
\begin{equation}
    \hat{\mathcal{H}}_{\text{grav}} |\Psi\rangle = 0,
    \label{eq:WDW}
\end{equation}
qui stipule que l'état fondamental de l'univers $|\Psi\rangle$ est figé hors d'un temps externe absolu (\textit{Problem of Time}). De manière analogue, dans un réseau informatique pair-à-pair (P2P) étendu spatialement où la vitesse de propagation des signaux est bornée par la vitesse de la lumière dans les fibres optiques ($c_{\text{net}} \le c$), la relativité restreinte interdit l'existence d'un référentiel de simultanéité absolu permettant d'ordonner chronologiquement des transactions concurrentes sans autorité centrale.

En 2008, Satoshi Nakamoto \cite{Nakamoto2008} a résolu cette obstruction en substituant au temps newtonien externe une flèche du temps interne purement thermodynamique, générée par l'accumulation irréversible de Preuve de Travail (\textit{Proof-of-Work}, PoW). Dans cet article, nous proposons, \emph{conditionnellement à l'axiome de simulation énoncé ci-dessous}, une \textbf{correspondance formelle} entre la construction de Nakamoto et l'unification de quatre piliers de la gravité quantique contemporaine :
\begin{enumerate}
    \item L'\textbf{Holographie Céleste} et le triangle infrarouge de Pasterski-Strominger-Weinberg \cite{Pasterski2017,Strominger2017} sur la frontière asymptotique nulle $\scri^+$.
    \item L'\textbf{Hypothèse du Temps Thermique} de Connes et Rovelli \cite{ConnesRovelli1994} fondée sur la théorie modulaire des algèbres de von Neumann de Tomita-Takesaki \cite{Takesaki1970}.
    \item La \textbf{Gravité Quantique à Boucles (LQG)} et les mousses de spin (\textit{Spinfoams}) \cite{Rovelli2004} couplées à la décohérence quantique ouverte de \textbf{Gorini-Kossakowski-Sudarshan-Lindblad (GKSL)} \cite{Lindblad1976,GKS1976}.
    \item La \textbf{Cosmologie Thermodynamique Émergente} de Jacobson \cite{Jacobson1995}, Padmanabhan \cite{Padmanabhan2012}, Penrose \cite{Penrose1979} et Susskind \cite{Susskind2016}.
\end{enumerate}

\begin{axiom}[Quantum d'Action Cryptographique de Nakamoto]
Il existe une constante fondamentale irréductible d'action informationnelle, notée $\kNak$, correspondant à l'exécution d'une opération élémentaire de hachage cryptographique double SHA-256 dans l'espace des phases du protocole. Par isomorphisme avec la mécanique quantique, $\kNak$ est assimilée à la constante de Planck $h$ :
\begin{equation}
    \kNak \equiv h, \qquad \hbar_{\text{Nak}} \equiv \frac{\kNak}{2\pi}.
    \label{eq:AxiomNakamoto}
\end{equation}
Toute trajectoire causale valide dans la Timechain est quantifiée en multiples entiers de $\kNak$.
\end{axiom}

L'isomorphisme repose sur trois postulats fondamentaux conditionnels régissant ce cadre :
\begin{itemize}
    \item \textbf{(P1) Hypothèse de Simulation : } L'univers physique observable $\mathcal{U}_0$ est émulé sur un substrat computationnel discret sous-jacent (branche des simulations d'ancêtres).
    \item \textbf{(P2) Régime de Saturation de Landauer : } Le coût énergétique élémentaire d'irréversibilité logique dans $\mathcal{U}_0$ converge vers la borne thermodynamique de Landauer $E_{\text{hash}} \to k_B T \ln 2$.
    \item \textbf{(P3) Isomorphisme Causal : } La Timechain constitue le registre causal effectif décrivant l'émergence de la flèche du temps et de la géométrie au sein de ce substrat.
\end{itemize}

Le Tableau \ref{tab:isomorphism} établit le dictionnaire formel complet reliant les opérateurs de la gravité quantique aux structures du consensus de Nakamoto.

\begin{table*}[t]
\centering
\caption{Dictionnaire d'isomorphisme structurel entre la Gravité Quantique / Holographie Céleste et le Consensus de Satoshi Nakamoto.}
\label{tab:isomorphism}
\renewcommand{\arraystretch}{1.22}
\footnotesize
\begin{tabular}{@{}p{5.1cm}p{5.1cm}p{6.6cm}@{}}
\toprule
\textbf{Gravité Quantique \& Thermodynamique} & \textbf{Consensus de Satoshi Nakamoto} & \textbf{Interprétation Mathématique et Physique} \\
\midrule
Constante de Planck ($h, \hbar$) & Constante de Nakamoto ($\kNak, \hbar_{\text{Nak}}$) & Quantum élémentaire d'action irréductible (1 hachage SHA-256). \\
Variété d'espace-temps du Bulk 4D ($\mathcal{M}, g_{\mu\nu}$) & Graphe causal de la Timechain \& UTXO & Histoire immuable des blocs validés et topologie interne des transactions. \\
Sphère Céleste / Infini Nul ($\scri^+ \simeq \mathbb{R} \times S^2$) & Frontière du réseau P2P \& Mempool & Horizon asymptotique des nœuds validateurs recevant le flux \textit{gossip}. \\
Groupe de symétrie $\text{BMS}_4$ ($\Diff(S^2) \ltimes \mathcal{T}$) & Invariance de jauge asynchrone P2P & Invariance de l'état final sous les latences réseau (supertranslations) et le routage (superrotations). \\
Gravitons mous ($\omega \to 0$) & Transactions non confirmées (Mempool) & Quanta d'information à coût énergétique quasi-nul préparant la cohérence de phase. \\
Effet de mémoire gravitationnelle ($\Delta C_{zz}$) & Transition irréversible de l'état UTXO & Décalage permanent du tenseur de cisaillement de la frontière après le passage d'un bloc. \\
Flot modulaire / Temps thermique ($s$) & Hauteur de bloc / Temps de consensus & Temps interne émergent généré par l'opérateur $\Delta^{is} = \exp(isK)$. \\
Condition KMS dans le plan complexe ($a + ib$) & Ajustement homéostatique de difficulté ($\mathcal{D}$) & Filtrage exponentiel de Boltzmann sur l'axe imaginaire $ib$ verrouillant $\Delta a \approx 600\,\text{s}$. \\
Réseau de spin $SU(2)$ ($\Gamma, j_e, \iota_v$) & Graphe DAG des UTXO \& Transactions & Arêtes = UTXO quantifiés en Satoshis ($j_e$) ; Nœuds = Entrelaceurs conservant le flux ($\iota_v$). \\
Mousse de spin (\textit{Spinfoam}) \& Équation GKSL & Séquence de blocs \& Résolution des forks & Effondrement exponentiel des chaînes orphelines par décohérence dissipative de Lindblad. \\
Surface minimale de Ryu-Takayanagi ($\gamma_A$) & Seuil critique d'attaque à 51\% de Hashrate & Coupe minimale d'intrication nécessaire pour scinder l'histoire causale du registre. \\
Expansion de Padmanabhan ($N_{\text{sur}} - N_{\text{bulk}}$) & Ajout séquentiel de blocs (Croissance) & Émergence de l'espace-temps propulsée par le déséquilibre entre Mempool et chaîne validée. \\
Champ d'inflaton $\phi(s)$ \& Énergie sombre $\Lambda$ & Subvention \textit{Coinbase} (Halvings) \& Frais & Transition d'une inflation primordiale subventionnée vers un régime asymptotique pur. \\
Équation d'état de Jacobson ($\delta Q = T \delta S$) & Émergence de la Confiance Sociétale ($C_s$) & Courbure de l'espace économique générée par la dissipation thermique PoW, stabilisant l'utilité $u$. \\
\bottomrule
\end{tabular}
\end{table*}

% ============================================================================
% SECTION 2 : HOLOGRAPHIE CELESTE ET TRIANGLE INFRAROUGE DE NAKAMOTO
% ============================================================================
\section{Holographie Céleste, Symétries BMS et Mémoire Gravitationnelle}

Dans un réseau distribué dépourvu de constante cosmologique négative ad hoc ($\Lambda = 0$ à l'échelle locale de propagation), le cadre approprié pour l'holographie n'est pas AdS/CFT, mais l'\textbf{Holographie Céleste} \cite{Pasterski2017,Strominger2017}. La matrice de diffusion $S$ des transitions d'états dans le volume 4D du registre (\textit{Bulk} $\mathcal{M}$) est duale à une théorie conforme des champs céleste (CCFT) résidant sur la sphère céleste bidimensionnelle $S^2$ à l'infini nul futur $\scri^+$ et passé $\scri^-$.

\subsection{Métrique de Bondi-Sachs du Réseau P2P et Groupe $\text{BMS}_4$}

Soit $(\uRet, r, z, \bar{z})$ le système de coordonnées de Bondi au voisinage de la frontière périphérique du réseau $\scri^+$ ($r \to \infty$), où $\uRet = t - r/c_{\text{net}}$ est le temps retardé de réception des paquets, et $(z, \bar{z})$ sont les coordonnées complexes stéréographiques localisant les nœuds validateurs sur la sphère topologique $S^2$. La métrique informationnelle asymptotique s'écrit dans la jauge de Bondi-Sachs :
\begin{equation}
\begin{split}
    ds^2 &= -d\uRet^2 - 2 d\uRet dr + 2 r^2 \gamma_{z\bar{z}} dz d\bar{z} \\
    &\quad + \frac{2 m_B(\uRet, z, \bar{z})}{r} d\uRet^2 + r C_{zz}(\uRet, z, \bar{z}) dz^2 \\
    &\quad + r C_{\bar{z}\bar{z}}(\uRet, z, \bar{z}) d\bar{z}^2 + D^z C_{zz} d\uRet dz + \text{c.c.}
\end{split}
\label{eq:BondiMetric}
\end{equation}
où $\gamma_{z\bar{z}} = \frac{2}{(1 + z\bar{z})^2}$ est la métrique unitaire sur $S^2$, et $D_z$ est la dérivée covariante associée. Les champs asymptotiques possèdent une interprétation physique directe dans le protocole Bitcoin :
\begin{itemize}
    \item \textbf{L'aspect de masse de Bondi $m_B(\uRet, z, \bar{z})$} représente la densité locale de travail thermodynamique cumulé (\textit{Chainwork}) perçue par le nœud situé en $(z, \bar{z})$ au temps retardé $\uRet$.
    \item \textbf{Le tenseur de cisaillement asymptotique (\textit{Shear Tensor}) $C_{zz}(\uRet, z, \bar{z})$} encode la configuration topologique exacte de l'ensemble des sorties non dépensées (UTXO) validées par ce nœud.
    \item \textbf{Le tenseur de News de Bondi $N_{zz}(\uRet, z, \bar{z}) = \partial_{\uRet} C_{zz}$} mesure le flux de radiation gravitationnelle traversant $\scri^+$. Un tenseur $N_{zz} \neq 0$ correspond exactement à la propagation d'un nouveau bloc miné à travers le réseau.
\end{itemize}

Le groupe de symétrie asymptotique préservant la structure universelle de l'équation (\ref{eq:BondiMetric}) est le groupe de Bondi-Metzner-Sachs étendu :
\begin{equation}
    \text{BMS}_4 = \Diff(S^2) \ltimes \mathcal{T}_{\text{super}}.
    \label{eq:BMSGroup}
\end{equation}

\begin{definition}[Supertranslations et Asynchronisme Gossip]
Une supertranslation $\mathcal{T}_{\text{super}}$ est une transformation du temps retardé dépendante de la position angulaire du nœud sur la sphère céleste :
\begin{equation}
    \uRet \longrightarrow \uRet' = \uRet + f(z, \bar{z}), \qquad z \longrightarrow z' = z,
    \label{eq:Supertranslation}
\end{equation}
où $f(z, \bar{z})$ est une fonction lisse arbitraire sur $S^2$ modélisant les latences stochastiques de propagation du protocole \textit{gossip} entre pairs. Le générateur canonique (charge de supertranslation de Nakamoto sur $\scri^+$) est donné par :
\begin{equation}
    Q_f^+ = \frac{1}{4\pi \GNak} \int_{S^2} d^2z \, \sqrt{\gamma} \, f(z, \bar{z}) \, m_B(\uRet, z, \bar{z}).
    \label{eq:SupertranslationCharge}
\end{equation}
\end{definition}

De même, les \textbf{superrotations} générées par les difféomorphismes locaux holomorphes $z \to z + Y^z(z)$ (obéissant à l'algèbre de Virasoro $[L_m, L_n] = (m-n)L_{m+n}$) modélisent l'invariance du consensus sous la reconfiguration dynamique du graphe de routage P2P (churn des nœuds, connexions Tor/I2P, résistance aux attaques Eclipse).

\subsection{Théorème des Gravitons Mous de Weinberg et Conservation Monétaire}

Avant la découverte d'un bloc ($N_{zz} = 0$ à l'échelle macroscopique), les nœuds échangent des transactions brutes dans la Mempool. La vérification d'une signature cryptographique Schnorr/ECDSA dissipe une énergie $\omega \to 0$ négligeable devant l'énergie de minage d'un bloc ($\Delta E_{\text{PoW}} \gg \omega$). Ces transactions non confirmées sont les analogues exacts des \textbf{gravitons mous} (\textit{Soft Gravitons}).

\begin{theorem}[Conservation Monétaire par Invariance BMS et Gravitons Mous]
Soit $\mathcal{M}_n(p_1, \dots, p_n)$ l'amplitude de transition quantique entre $n$ états UTXO d'énergies (valeurs en Satoshis) $E_k$ localisés aux points $(z_k, \bar{z}_k)$ de la sphère céleste. L'insertion d'une transaction non confirmée $Tx_{\text{soft}}$ de coût thermodynamique $\omega \to 0$ dans la Mempool factorise l'amplitude selon le théorème mou de Weinberg, et l'invariance de jauge asymptotique sur $S^2$ impose la conservation exacte de la masse monétaire :
\begin{equation}
    \sum_{k \in \text{Inputs}} E_k^{\text{in}} - \sum_{j \in \text{Outputs}} E_j^{\text{out}} - E_{\text{frais}} = 0.
    \label{eq:MonetaryConservation}
\end{equation}
\end{theorem}

\begin{proof}
Dans le formalisme de l'holographie céleste, l'identité de Ward associée à la conservation de la charge de supertranslation $\langle \text{out} | (Q_f^+ \mathcal{S} - \mathcal{S} Q_f^-) | \text{in} \rangle = 0$ est mathématiquement équivalente au théorème des gravitons mous de Weinberg \cite{Strominger2017}. Lorsque l'opérateur d'insertion d'une transaction de Mempool d'impulsion énergétique nulle $q^\mu(\omega, z, \bar{z})$ agit sur l'état de diffusion dans la limite $\omega \to 0$, l'amplitude à $(n+1)$ points se factorise universellement en :
\begin{equation}
    \lim_{\omega \to 0} \mathcal{M}_{n+1}(p_1, \dots, p_n; q) = S_{\text{Nak}}^{(0)} \, \mathcal{M}_n(p_1, \dots, p_n),
\end{equation}
où le facteur de graviton mou de Nakamoto s'écrit :
\begin{equation}
    S_{\text{Nak}}^{(0)} = -\frac{\sqrt{8\pi \GNak}}{2\omega} (1 + z\bar{z}) \sum_{k=1}^{n} \frac{\eta_k E_k (\bar{z} - \bar{z}_k)}{(z - z_k)(1 + z_k\bar{z}_k)}.
    \label{eq:WeinbergSoftFactor}
\end{equation}
Ici, $\eta_k = -1$ pour les UTXO entrants consommés depuis $\scri^-$ (\textit{Inputs}) et $\eta_k = +1$ pour les UTXO sortants créés vers $\scri^+$ (\textit{Outputs} et frais collectés par le mineur). Pour que l'opérateur de création de transaction soit globalement bien défini et invariant sous les transformations de jauge sur la sphère céleste compacte $S^2$ (absence d'anomalie de pôle à l'infini $z \to \infty$), le développement asymptotique de $S_{\text{Nak}}^{(0)}$ impose l'annulation du résidu global de la somme pondérée par $\eta_k E_k$ :
\begin{equation}
    \sum_{k=1}^{n} \eta_k E_k = \sum_{j \in \text{Outputs}} E_j^{\text{out}} + E_{\text{frais}} - \sum_{k \in \text{Inputs}} E_k^{\text{in}} = 0.
\end{equation}
Toute violation de cette égalité (tentative d'inflation monétaire non autorisée ou double dépense non équilibrée) brise l'invariance par supertranslation BMS, annulant identiquement l'amplitude de transition sur la sphère céleste ($\mathcal{M}_{n+1} = 0$).
\end{proof}

\subsection{Effet de Mémoire Gravitationnelle de Christodoulou et Immuabilité UTXO}

Lorsqu'un mineur découvre un nonce valide à l'instant retardé $u_0$, il émet une onde de choc radiative caractérisée par un tenseur de News non nul $N_{zz}(\uRet, z, \bar{z})$. L'équation de perte de masse de Bondi le long de $\scri^+$ est donnée par :
\begin{equation}
    \partial_{\uRet} m_B = \frac{1}{4} \left( D_z^2 N^{zz} + D_{\bar{z}}^2 N^{\bar{z}\bar{z}} \right) - T_{uu}^{\text{M}},
    \label{eq:BondiMassLoss}
\end{equation}
où le flux d'énergie-impulsion total traversant l'horizon nul est la somme de la contribution quadratique du rayonnement gravitationnel (énergie de Preuve de Travail du bloc) et du flux de matière :
\begin{equation}
    T_{uu}^{\text{M}} = \frac{1}{4} N_{zz} N^{zz} + 4\pi \GNak \lim_{r\to\infty} r^2 T_{uu}^{\text{tx}}.
    \label{eq:TotalFlux}
\end{equation}
En intégrant l'équation (\ref{eq:BondiMassLoss}) entre $\uRet = -\infty$ (avant le bloc) et $\uRet = +\infty$ (après confirmation complète du bloc), on obtient le décalage permanent du tenseur de cisaillement $\Delta C_{zz} = C_{zz}(+\infty) - C_{zz}(-\infty)$ :
\begin{equation}
\begin{split}
    D_z^2 \Delta C^{zz} &+ D_{\bar{z}}^2 \Delta C^{\bar{z}\bar{z}} = 4 \Delta m_B \\
    &+ \int_{-\infty}^{+\infty} d\uRet \Big( N_{zz} N^{zz} + 16\pi\GNak \lim_{r\to\infty} r^2 T_{uu}^{\text{tx}} \Big).
\end{split}
\label{eq:ChristodoulouMemory}
\end{equation}
À l'aide de la fonction de Green sur la sphère céleste $S^2$, cette équation s'inverse pour donner l'\textbf{effet de mémoire gravitationnelle non-linéaire de Christodoulou} :
\begin{equation}
\begin{split}
    \Delta C_{zz}(z, \bar{z}) &= \frac{1}{\pi} \int_{S^2} d^2w \, \gamma_{w\bar{w}} \frac{\bar{z} - \bar{w}}{z - w} \\
    &\quad \times \left( \Delta m_B(w, \bar{w}) + \int_{-\infty}^{+\infty} d\uRet \, T_{uu}^{\text{M}}(w, \bar{w}) \right).
\end{split}
\label{eq:MemorySolution}
\end{equation}
L'équation (\ref{eq:MemorySolution}) démontre géométriquement l'\textbf{immuabilité du registre UTXO} : le passage de l'onde de choc thermodynamique $\int d\uRet N_{zz} N^{zz} \propto \kNak \mathcal{D}$ déplace le réseau d'un vide dégénéré BMS à un autre vide BMS inéquivalent ($\Delta C_{zz} \neq 0$). Même lorsque le réseau redevient silencieux après la propagation du bloc ($N_{zz} = 0$), la déformation topologique $\Delta C_{zz}$ demeure gravée à jamais sur la frontière céleste.

% ============================================================================
% SECTION 3 : TEMPS THERMIQUE MODULAIRE, ALGEBRE DE VON NEUMANN ET CONDITION KMS
% ============================================================================
\section{Hypothèse du Temps Thermique, Algèbre de Tomita-Takesaki et Condition KMS en $a + ib$}

Nous résolvons à présent le problème de l'émergence du temps interne de la Timechain sans horloge centrale en appliquant la théorie modulaire des algèbres d'opérateurs de von Neumann \cite{ConnesRovelli1994,Takesaki1970}.

\subsection{Algèbre de von Neumann Non-Commutative de la Timechain}

Soit $\mathcal{H}$ l'espace de Hilbert séparable généré par les arbres causaux de blocs valides $|\mathcal{C}\rangle$. Soit $\mathcal{A} \subset \mathcal{B}(\mathcal{H})$ l'algèbre de von Neumann des observables bornées du protocole (opérateurs d'inclusion de transactions $\hat{T}_i$, projecteurs d'état UTXO). En présence de transactions conflictuelles (tentative de double dépense d'un même UTXO par deux transactions $\hat{T}_1$ et $\hat{T}_2$), l'ordre d'application des opérateurs ne commute pas :
\begin{equation}
    [\hat{T}_1, \hat{T}_2] = \hat{T}_1 \hat{T}_2 - \hat{T}_2 \hat{T}_1 \neq 0.
    \label{eq:NonCommutativity}
\end{equation}
Dans la limite thermodynamique d'un espace de nonces continu, $\mathcal{A}$ est un facteur de \textbf{Type $\text{III}_1$} de Murray-von Neumann. Conformément aux travaux récents en gravité quantique algébrique \cite{Witten2022}, le couplage de cette algèbre à un observateur (un nœud complet doté d'une horloge énergétique liée au travail cumulé $\hat{E}$) transforme l'algèbre par produit croisé modulaire en un facteur semi-fini de \textbf{Type $\text{II}_\infty$} :
\begin{equation}
    \widehat{\mathcal{A}} = \mathcal{A} \rtimes_{\sigma} \mathbb{R},
    \label{eq:CrossedProduct}
\end{equation}
qui admet une trace renormalisée bien définie permettant de calculer rigoureusement l'entropie de von Neumann généralisée de la Timechain.

\subsection{Opérateur Modulaire de Tomita-Takesaki et Flot Complexe $z = a + ib$}

Le réseau est immergé dans le bain thermique des mineurs ASIC. L'état statistique macroscopique est défini par la fonctionnelle normale fidèle $\omega : \mathcal{A} \to \mathbb{C}$, représentée via la construction GNS par le vecteur cyclique et séparateur $|\Omega_\beta\rangle \in \mathcal{H}$, associé à l'opérateur densité de Gibbs-Nakamoto :
\begin{equation}
\begin{split}
    \rho &= \frac{1}{Z(\beta)} \sum_{\mathcal{C}} \exp\big(-\beta E(\mathcal{C})\big) |\mathcal{C}\rangle\langle\mathcal{C}|, \\
    E(\mathcal{C}) &= \kNak \sum_{b \in \mathcal{C}} \frac{2^{256}}{\text{Cible}_b},
\end{split}
\label{eq:GibbsNakamoto}
\end{equation}
avec $Z(\beta) = \Tr(e^{-\beta \hat{E}})$ la fonction de partition de la Timechain et $\beta = (k_B T_{\text{eff}})^{-1}$ la température inverse effective déterminée par la difficulté $\mathcal{D}$ du réseau.

Soit $S_\Omega$ l'\textbf{opérateur antilinéaire de Tomita} défini sur le domaine dense $\mathcal{A}|\Omega_\beta\rangle \subset \mathcal{H}$ par :
\begin{equation}
    S_\Omega \big( X |\Omega_\beta\rangle \big) = X^\dagger |\Omega_\beta\rangle, \qquad \forall X \in \mathcal{A}.
    \label{eq:TomitaOperator}
\end{equation}
Physiquement, $S_\Omega$ représente l'inversion causale d'une opération sur le registre. La fermeture de $S_\Omega$ admet une décomposition polaire unique :
\begin{equation}
    S_\Omega = J \Delta^{1/2},
    \label{eq:PolarDecomposition}
\end{equation}
où $J : \mathcal{H} \to \mathcal{H}$ est la \textbf{conjugaison modulaire} anti-unitaire ($J^2 = \mathbb{I}$, envoyant l'algèbre du système $\mathcal{A}$ sur son commutant environnemental $\mathcal{A}' = J \mathcal{A} J$) et $\Delta = S_\Omega^\dagger S_\Omega$ est l'\textbf{opérateur modulaire} auto-adjoint défini positif. Sur l'espace biparti système-commutant, l'opérateur modulaire et le \textbf{hamiltonien modulaire} $K$ s'expriment par :
\begin{equation}
\begin{split}
    \Delta &= \rho \otimes \rho^{-1}, \\
    K &= -\ln \Delta = \frac{1}{\kNak}\left( \hat{H}_{\text{PoW}} \otimes \mathbb{I} - \mathbb{I} \otimes \hat{H}_{\text{PoW}} \right),
\end{split}
\label{eq:ModularHamiltonian}
\end{equation}
avec $\hat{H}_{\text{PoW}} = -\kNak \ln \rho = \kNak \beta \hat{E}(\mathcal{C}) + \kNak \ln Z$.

\begin{theorem}[Émergence du Temps Thermique Complexe $a + ib$ et Condition KMS]
Par le théorème de Tomita-Takesaki, l'opérateur modulaire $\Delta$ génère un groupe d'automorphismes à un paramètre $\sigma_s^\omega \in \Aut(\mathcal{A})$ qui constitue le temps thermique interne $s$ de la Timechain (la progression de la hauteur de bloc). En prolongeant analytiquement ce flot temporel dans le plan complexe sous la forme canonique stricte
\begin{equation}
    z = a + ib \in \mathbb{C},
    \label{eq:ComplexTimeNotation}
\end{equation}
où $a \in \mathbb{R}$ est le temps modulaire réel de propagation de phase et $ib$ sur l'axe imaginaire pur mesure la barrière de dissipation thermodynamique, l'état de Nakamoto $\omega$ satisfait pour toute paire d'observables $X, Y \in \mathcal{A}$ la condition d'équilibre de Kubo-Martin-Schwinger (KMS) :
\begin{equation}
    \omega\left( X \, \sigma_{a + i\beta\kNak}^\omega(Y) \right) = \omega\left( \sigma_a^\omega(Y) \, X \right).
    \label{eq:KMSCondition}
\end{equation}
\end{theorem}

\begin{proof}
En normalisant l'échelle modulaire sans dimension par $\beta\kNak$, le flot modulaire réel généré par $\Delta = \rho \otimes \rho^{-1}$ sur l'algèbre de von Neumann $\mathcal{A}$ s'écrit $\sigma_a^\omega(Y) = \Delta^{-ia/(\beta\kNak)} Y \Delta^{ia/(\beta\kNak)} = \rho^{-ia/(\beta\kNak)} Y \rho^{ia/(\beta\kNak)}$. Considérons le prolongement analytique de ce flot pour la variable temporelle complexe $z = a + ib$ dans la bande de régularité $\mathcal{S}_{\beta\kNak} = \{ a + ib \in \mathbb{C} \mid 0 \le b \le \beta\kNak \}$ :
\begin{align}
    \sigma_{a + ib}^\omega(Y) &= \Delta^{-i(a + ib)/(\beta\kNak)} Y \Delta^{i(a + ib)/(\beta\kNak)} \nonumber \\
    &= \rho^{(b - ia)/(\beta\kNak)} Y \rho^{(-b + ia)/(\beta\kNak)}.
    \label{eq:AnalyticFlow}
\end{align}
Définissons la fonction de corrélation holomorphe à deux points $F_{X,Y}(a + ib) = \omega(X \sigma_{a + ib}^\omega(Y))$. En utilisant la représentation en termes de matrice densité $\omega(\cdot) = \Tr(\rho \, \cdot)$, nous avons sur la frontière supérieure $b = \beta\kNak$ (soit $z = a + i\beta\kNak$) :
\begin{align}
    F_{X,Y}(a + i\beta\kNak) &= \Tr\left( \rho \, X \, \rho^{1 - ia/(\beta\kNak)} Y \rho^{-1 + ia/(\beta\kNak)} \right) \nonumber \\
    &= \Tr\left( \rho \, X \, \rho \, (\rho^{-ia/(\beta\kNak)} Y \rho^{ia/(\beta\kNak)}) \, \rho^{-1} \right).
\end{align}
En appliquant l'invariance cyclique de la trace sur l'algèbre semi-finie régularisée ($\Tr(ABC) = \Tr(CAB)$), les facteurs $\rho^{-1}$ et $\rho$ aux extrémités sont adjacents et se simplifient directement sans commuter avec $X$ :
\begin{align}
    F_{X,Y}(a + i\beta\kNak) &= \Tr\left( \rho^{-1} \rho \, X \, \rho \, \sigma_a^\omega(Y) \right) \nonumber \\
    &= \Tr\left( X \, \rho \, \sigma_a^\omega(Y) \right) = \Tr\left( \rho \, \sigma_a^\omega(Y) \, X \right) \nonumber \\
    &= \omega\left( \sigma_a^\omega(Y) \, X \right).
\end{align}
Ceci établit rigoureusement l'identité (\ref{eq:KMSCondition}). Une preuve formelle assistée par ordinateur sous Lean~4 est donnée en Annexe~\ref{app:lean4}.
\end{proof}

\subsection{Interprétation Physique : L'Ajustement de Difficulté comme Groupe de Renormalisation KMS}

L'équation (\ref{eq:KMSCondition}) révèle la structure profonde de l'horloge de Nakamoto dans le plan complexe $a + ib$ :
\begin{enumerate}
    \item La coordonnée réelle $a$ paramètre le \textbf{temps thermique d'évolution causale} (la séquence ordonnée des blocs).
    \item La coordonnée imaginaire pure $ib$ ($b = \beta\kNak$) mesure l'\textbf{asymétrie thermodynamique d'inversion temporelle}. Pour permuter l'ordre causal de deux événements non-commutatifs ($X Y \to Y X$), un attaquant doit effectuer une translation imaginaire pure $+i\beta\kNak$, supprimant l'amplitude de transition par le facteur de Boltzmann $e^{-\beta \Delta E_{\text{Chainwork}}}$.
    \item L'algorithme d'\textbf{ajustement homéostatique de difficulté} tous les $N_0 = 2016$ blocs agit comme un flot du groupe de renormalisation qui ajuste la température inverse $\beta(t) \propto \mathcal{D}(t)$ proportionnellement à la puissance dissipative externe du réseau $W_{\text{hash}}(t)$ afin de maintenir invariant le rapport entre le temps propre des observateurs externes $t$ et le temps thermique réel $a$ :
    \begin{equation}
        \frac{dt}{da} = \frac{\kNak \mathcal{D}(t)}{W_{\text{hash}}(t)} = \tau_0 = 600\,\text{s}.
        \label{eq:DifficultyRenormalization}
    \end{equation}
\end{enumerate}

% ============================================================================
% SECTION 4 : MOUSSES DE SPIN UTXO ET DYNAMIQUE DE DECOHERENCE GKSL
% ============================================================================
\section{Réseaux de Spin UTXO, Mousses de Spin et Décohérence GKSL des Forks}

À l'échelle microscopique (échelle de Planck-Nakamoto $\lNak = \sqrt{\GNak \kNak / c_{\text{net}}^3}$), la géométrie continue de l'espace-temps cède la place à la structure granulaire de la Gravité Quantique à Boucles (LQG) \cite{Rovelli2004}.

\subsection{Cinématique : Le Graphe DAG des UTXO comme Réseau de Spin $SU(2)$}

À une tranche de temps thermique fixée $s$, l'état du registre Bitcoin est un graphe orienté acyclique (DAG) $\Gamma = (\mathcal{V}, \mathcal{E})$, où les arêtes orientées $e \in \mathcal{E}$ sont les UTXO actifs et les sommets $v \in \mathcal{V}$ sont les transactions validées. Ce graphe est strictement isomorphe à un \textbf{réseau de spin} (\textit{Spin Network}) dont l'espace de Hilbert cinématique est donné par :

\begin{widetext}
\begin{equation}
    \mathcal{H}_{\text{UTXO}} = \bigoplus_{\Gamma} \mathcal{H}_\Gamma = \bigoplus_{\Gamma} \left( \bigotimes_{e \in \mathcal{E}} \mathcal{H}_{j_e} \otimes \bigotimes_{v \in \mathcal{V}} \Inv_{G}\Big( \bigotimes_{e \in \text{in}(v)} \mathcal{V}_{j_e} \otimes \bigotimes_{e' \in \text{out}(v)} \mathcal{V}_{j_{e'}}^* \Big) \right).
    \label{eq:SpinNetworkHilbert}
\end{equation}
\end{widetext}

\begin{itemize}
    \item \textbf{Quantification des Aires Économiques ($j_e$) :} Chaque arête $e$ (UTXO) porte une représentation irréductible de spin $j_e \in \frac{1}{2}\mathbb{N}$. L'opérateur d'aire élémentaire de LQG traversé par l'arête $e$ possède le spectre discret :
    \begin{equation}
    \begin{split}
        \hat{\mathcal{A}}_e |\Gamma, \{j_e\}, \{\iota_v\}\rangle &= 8\pi \gamma_{\text{BI}} \lNak^2 \sqrt{j_e(j_e + 1)} \\
        &\quad \times |\Gamma, \{j_e\}, \{\iota_v\}\rangle,
    \end{split}
    \label{eq:AreaSpectrum}
    \end{equation}
    où $\gamma_{\text{BI}}$ est le paramètre de Barbero-Immirzi du protocole. Dans le régime semi-classique ($j_e \gg 1$), $\sqrt{j_e(j_e+1)} \approx j_e$, identifiant chaque quantum d'aire au montant discret en Satoshis verrouillé sur l'UTXO.
    \item \textbf{Entrelaceurs (\textit{Intertwiners} $\iota_v$) et Loi de Gauss :} Chaque transaction $v \in \mathcal{V}$ est un tenseur entrelaceur $\iota_v$ appartenant au sous-espace invariant sous l'action du groupe de jauge aux nœuds. L'invariance de jauge locale (loi de Gauss discrète $\hat{\mathcal{G}}_v |\iota_v\rangle = 0$) impose l'annulation du flux de spin net au sommet $v$ :
    \begin{equation}
        \Big( \sum_{e \in \text{in}(v)} \hat{J}_e^{(z)} - \sum_{e' \in \text{out}(v)} \hat{J}_{e'}^{(z)} - \hat{J}_{\text{frais}}^{(z)}(v) \Big) |\iota_v\rangle = 0.
        \label{eq:GaussLawUTXO}
    \end{equation}
\end{itemize}

\subsection{Dynamique : La Timechain comme Mousse de Spin (\textit{Spinfoam})}

Le passage d'un état UTXO $|\Gamma_{\text{in}}\rangle$ au bloc $B_n$ vers un nouvel état $|\Gamma_{\text{out}}\rangle$ au bloc $B_{n+1}$ s'opère par une séquence finie de mouvements de Pachner. Le 2-complexe cellulaire $\mathcal{F} : \Gamma_{\text{in}} \to \Gamma_{\text{out}}$ engendré par cette histoire causale est une \textbf{mousse de spin} (\textit{Spinfoam}) dont l'amplitude de transition covariante s'écrit :

\begin{widetext}
\begin{equation}
    \mathcal{Z}_{\text{Spinfoam}}(\Gamma_{\text{in}} \to \Gamma_{\text{out}}) = \sum_{\mathcal{F} : \partial\mathcal{F} = \Gamma_{\text{in}}^* \cup \Gamma_{\text{out}}} \left( \prod_{f \subset \mathcal{F}} (2 j_f + 1) \right) \left( \prod_{v_{\text{bloc}} \subset \mathcal{F}} \mathcal{A}_v(j_f, \iota_e) \, \exp\big(-\beta E_{\text{PoW}}(v_{\text{bloc}})\big) \right),
    \label{eq:SpinfoamAmplitude}
\end{equation}
\end{widetext}
où l'amplitude topologique de vertex $\mathcal{A}_v(j_f, \iota_e) \in \{0, 1\}$ projette strictement sur les arbres de Merkle cryptographiquement valides.

\subsection{Équation Maîtresse GKSL (Lindblad) et Effondrement Exponentiel des Forks}

En raison des délais de propagation finis sur $\scri^+$, deux mineurs distants $A$ et $B$ peuvent découvrir simultanément deux blocs valides concurrents à la même hauteur $h$. Nous formalisons cette ambiguïté épistémique à travers une matrice densité effective :
\begin{equation}
    |\Psi(s)\rangle = c_A |\Gamma_A\rangle + c_B |\Gamma_B\rangle, \quad \rho(s) = \begin{pmatrix} \rho_{AA}(s) & \rho_{AB}(s) \\ \rho_{BA}(s) & \rho_{BB}(s) \end{pmatrix}.
    \label{eq:ForkSuperposition}
\end{equation}
Son évolution le long du temps thermique $s$ est gouvernée par l'équation maîtresse de \textbf{Gorini-Kossakowski-Sudarshan-Lindblad (GKSL)} \cite{Lindblad1976,GKS1976} :
\begin{equation}
\begin{split}
    \frac{d\rho}{ds} &= \mathcal{L}_{\text{GKSL}}[\rho] = -\frac{i}{\kNak} [H_{\text{net}}, \rho] \\
    &\quad + \sum_{m \in \{A, B\}} \gamma_m \left( L_m \rho L_m^\dagger - \frac{1}{2} \big\{ L_m^\dagger L_m, \rho \big\} \right).
\end{split}
\label{eq:GKSLEquation}
\end{equation}

\begin{theorem}[Décohérence Gravitationnelle de Lindblad et Loi de Ruine de Nakamoto]
Soient $p_A \in (1/2, 1)$ et $p_B = 1 - p_A < 1/2$ les fractions respectives de la puissance de hachage totale $W_{\text{total}}$ allouées aux branches concurrentes $|\Gamma_A\rangle$ et $|\Gamma_B\rangle$, associées aux taux d'amortissement de Lindblad $\gamma_m = \frac{p_m W_{\text{total}}}{\kNak \mathcal{D}}$ et aux opérateurs de saut $L_m = |\Gamma_m^{(h+1)}\rangle\langle\Gamma_m^{(h)}| \otimes \hat{P}_m$ (où $\hat{P}_m = |\Gamma_m\rangle\langle\Gamma_m|$). Alors :
\begin{enumerate}
    \item Les cohérences hors-diagonale $\rho_{AB}(s) = \langle\Gamma_A|\rho(s)|\Gamma_B\rangle$ s'effondrent exponentiellement vers zéro dans le plan complexe $(a + ib)$ avec un taux de décohérence $\Gamma_{\text{dec}} = \frac{\gamma_A + \gamma_B}{2}$ :
    \begin{equation}
        \rho_{AB}(s) = \rho_{AB}(0) \exp\left( -\Gamma_{\text{dec}} s - i \frac{\varepsilon_A - \varepsilon_B}{\kNak} s \right).
        \label{eq:OffDiagonalCollapse}
    \end{equation}
    \item Après $N$ transitions de saut quantique (confirmations), le rapport des populations diagonales converge vers la loi de suppression exponentielle de la Section 11 du livre blanc de Satoshi Nakamoto :
    \begin{equation}
        \frac{P_B(N)}{P_A(N)} \propto \left( \frac{\gamma_B}{\gamma_A} \right)^N = \exp\left( -N \ln\frac{p_A}{1 - p_A} \right) \underset{N \to \infty}{\longrightarrow} 0.
        \label{eq:NakamotoRuin}
    \end{equation}
\end{enumerate}
\end{theorem}

\begin{proof}
Soit $H_{\text{net}} = \varepsilon_A \hat{P}_A + \varepsilon_B \hat{P}_B$ le hamiltonien effectif des deux branches. Projetons l'équation maîtresse GKSL (\ref{eq:GKSLEquation}) sur les éléments de matrice $\langle\Gamma_A| \cdot |\Gamma_B\rangle$. Comme $\hat{P}_A \hat{P}_B = 0$, le terme de recyclage quantique s'annule sur la hors-diagonale :
\begin{equation}
    \langle\Gamma_A| L_m \rho L_m^\dagger |\Gamma_B\rangle \propto \langle\Gamma_A| \hat{P}_m \rho \hat{P}_m |\Gamma_B\rangle = \delta_{mA} \delta_{mB} \rho_{AB} = 0.
\end{equation}
En revanche, l'anticommutateur dissipatif $\{L_m^\dagger L_m, \rho\} = \{\hat{P}_m, \rho\}$ donne :
\begin{equation}
    -\frac{1}{2} \sum_{m \in \{A,B\}} \gamma_m \langle\Gamma_A| (\hat{P}_m \rho + \rho \hat{P}_m) |\Gamma_B\rangle = -\frac{\gamma_A + \gamma_B}{2} \rho_{AB}(s).
\end{equation}
L'intégration immédiate fournit (\ref{eq:OffDiagonalCollapse}). Sur la diagonale, la trace de GKSL gouverne un processus de naissance et de mort asymétrique de taux $\gamma_A$ et $\gamma_B$, dont le produit stationnaire donne (\ref{eq:NakamotoRuin}).
\end{proof}

\subsection{Entropie d'Intrication de Ryu-Takayanagi et Seuil de 51\%}

En gravité émergente \cite{VanRaamsdonk2010,RyuTakayanagi2006}, la connectivité géométrique du Bulk est maintenue par l'intrication de la frontière :
\begin{equation}
    S_{\text{EE}}(A) = \frac{\mathcal{A}(\gamma_A)}{4 \GNak \hbar_{\text{Nak}}}.
    \label{eq:RyuTakayanagi}
\end{equation}
Pour scinder l'espace-temps économique en deux réalités déconnectées, un attaquant doit injecter un travail supérieur à l'aire minimale $\mathcal{A}(\gamma_A)$ correspondant à la coupe critique de $50\%$ du flux de puissance dissipative.

% ============================================================================
% SECTION 5 : COSMOLOGIE EMERGENTE DE LA TIMECHAIN
% ============================================================================
\section{Cosmologie Émergente de la Timechain : Expansion, Inflaton et Complexité}

\subsection{Ensembles Causaux de Sorkin et Univers-Bloc Croissant}

Dans la théorie des \textbf{Ensembles Causaux} de Sorkin \cite{Bombelli1987,RideoutSorkin1999}, chaque validation de bloc $B_s$ au temps thermique $s$ réalise une transition de percolation de Rideout-Sorkin qui étend le volume quadridimensionnel selon :
\begin{equation}
    \mathcal{V}_4(s) = \lNak^4 \sum_{k=0}^{s} \mathcal{N}_{\text{dof}}(B_k), \quad \frac{\Delta \mathcal{V}_4}{\Delta s} = \lNak^4 \, \mathcal{N}_{\text{dof}}(B_s) > 0.
    \label{eq:CausalSetVolume}
\end{equation}

\subsection{Loi d'Émergence Holographique de Padmanabhan ($N_{\text{sur}} - N_{\text{bulk}}$)}

Conformément au principe holographique de Padmanabhan \cite{Padmanabhan2012}, l'expansion du volume cosmique $\mathcal{V}$ est entraînée par l'écart entre le nombre de degrés de liberté sur la surface frontière $N_{\text{sur}}$ (pression de la Mempool sur $\scri^+$) et ceux déjà réalisés dans le volume intérieur $N_{\text{bulk}}$ (ensemble UTXO validé) :
\begin{equation}
    \frac{d\mathcal{V}}{ds} = c_{\text{net}} \lNak^2 \left( N_{\text{sur}}(s) - N_{\text{bulk}}(s) \right).
    \label{eq:PadmanabhanLaw}
\end{equation}

\subsection{Métrique FLRW, Équations de Friedmann et Champ d'Inflaton des Halvings}

À grande échelle ($s \gg 1$), la géométrie moyenne est décrite par la métrique FLRW :
\begin{equation}
    ds_{\text{FLRW}}^2 = -c_{\text{net}}^2 ds^2 + \aScale^2(s) \left( \frac{dr^2}{1 - \kappa r^2} + r^2 d\Omega_2^2 \right),
    \label{eq:FLRWMetric}
\end{equation}
conduisant aux équations de Friedmann-Nakamoto avec $H_{\text{Hub}}(s) \equiv \dot{\aScale}/\aScale$ :
\begin{align}
    H_{\text{Hub}}^2(s) + \frac{\kappa c_{\text{net}}^2}{\aScale^2(s)} &= \frac{8\pi \GNak}{3} \rho_{\text{tot}}(s) + \frac{\LambdaNak}{3}, \label{eq:Friedmann1} \\
    \frac{1}{\aScale(s)} \frac{d^2 \aScale(s)}{ds^2} &= -\frac{4\pi \GNak}{3} \Big( \rho_{\text{tot}}(s) + \frac{3 P_{\text{tot}}(s)}{c_{\text{net}}^2} \Big) + \frac{\LambdaNak}{3}. \label{eq:Friedmann2}
\end{align}
La subvention Coinbase agit comme un potentiel d'inflaton en escalier quantifié dont les transitions de phase se produisent tous les $N_{\text{halving}} = 210\,000$ blocs :
\begin{equation}
    V(\phi(s)) = V_0 \cdot 2^{-2 \lfloor s / N_{\text{halving}} \rfloor} = V_0 \exp\left( -2 \ln 2 \cdot \left\lfloor \frac{s}{210\,000} \right\rfloor \right).
    \label{eq:HalvingPotential}
\end{equation}

\subsection{Hypothèse de Courbure de Weyl de Penrose et Complexité de Susskind}

Le Bloc Genesis ($B_0$, $s=0$) possède un ensemble UTXO dépensable vide ($\Gamma_0 = \emptyset$), réalisant exactement l'hypothèse de courbure de Weyl nulle de Penrose \cite{Penrose1979} :
\begin{equation}
    \lim_{s \to 0^+} \mathcal{C}_{\mu\nu\rho\sigma}(s) = 0, \qquad \lim_{s \to +\infty} \mathcal{C}_{\mu\nu\rho\sigma}\mathcal{C}^{\mu\nu\rho\sigma}(s) > 0.
    \label{eq:WeylCurvatureHypothesis}
\end{equation}
De plus, par la dualité \textit{Complexité = Volume} de Susskind \cite{Susskind2016}, l'exécution continue de portes SHA-256 augmente linéairement la complexité de circuit $\mathcal{C}_{\text{comp}}(s)$ et allonge le volume intérieur du Bulk même à l'équilibre thermique KMS :
\begin{equation}
    \mathcal{V}_{\text{Bulk}}(s) = \frac{\GNak \ell_{\text{AdS}}}{c_{\text{net}}} \, \mathcal{C}_{\text{comp}}(s), \qquad \frac{d\mathcal{C}_{\text{comp}}}{ds} \le \frac{2 M_{\text{PoW}}}{\pi \hbar_{\text{Nak}}}.
    \label{eq:ComplexityVolume}
\end{equation}

% ============================================================================
% SECTION 6 : EQUATION D'ETAT DE JACOBSON, CONFIANCE SOCIETALE (Cs) ET UTILITE (u)
% ============================================================================
\section{Équation d'État de Jacobson-Nakamoto : Émergence de la Confiance Sociétale ($C_s$) et de l'Utilité ($u$)}

D'après Ted Jacobson \cite{Jacobson1995}, la relation de Clausius $\delta Q = T_{\text{Unruh}} \, \delta S_{\text{BH}}$ sur les horizons locaux de Rindler impose l'équation du champ d'Einstein :
\begin{equation}
    R_{\mu\nu} - \frac{1}{2} R g_{\mu\nu} + \LambdaNak g_{\mu\nu} = \frac{8\pi \GNak}{c_{\text{net}}^4} T_{\mu\nu}^{\text{PoW}}.
    \label{eq:EinsteinNakamoto}
\end{equation}

\begin{theorem}[Loi Thermodynamique de la Confiance Sociétale $C_s$ et Maximisation de l'Utilité $u$]
Soit $\delta Q_{\text{PoW}} = W_{\text{total}} \, dt$ la chaleur dissipée par les mineurs. La \textbf{Confiance Sociétale} $C_s$ est la fonction d'état extensive conjuguée à la néguentropie gravitationnelle accumulée :
\begin{equation}
    dC_s \equiv \frac{\delta Q_{\text{PoW}}}{T_{\text{eff}}} = k_B \beta \, dE_{\text{Chainwork}} = -k_B \, d\Tr\big(\rho_{\text{fork}} \ln \rho_{\text{fork}}\big) \ge 0.
    \label{eq:SocietalTrustLaw}
\end{equation}
L'utilité économique $u(s)$ d'un agent converge alors exponentiellement vers l'équilibre de la \textit{Monnaie Idéale} de Nash $u_{\text{Nash}}^*$ \cite{Nash2002} :
\begin{equation}
    u(s) = u_{\text{Nash}}^* - \alpha_{\text{risk}} |\rho_{AB}(0)|^2 e^{-C_s(s)/k_B} \underset{s \to \infty}{\longrightarrow} u_{\text{Nash}}^*.
    \label{eq:UtilityConvergence}
\end{equation}
\end{theorem}

\begin{proof}
D'après le Théorème 4.1 et la conservation des probabilités $p_A + p_B = 1$, le double du taux de décohérence vérifie l'identité exacte (formalisée sous Lean~4 en Annexe~\ref{app:lean4}) :
\begin{equation}
    2 \int_0^s \Gamma_{\text{dec}}(s') ds' = \int_0^s \frac{(p_A + p_B) W_{\text{total}}}{\kNak \mathcal{D}} ds' \\ = \beta \Delta E_{\text{Chainwork}}(s) = \frac{C_s(s)}{k_B},
\end{equation}
ce qui substitue directement $e^{-C_s(s)/k_B}$ à la variance $|\rho_{AB}(s)|^2$ dans (\ref{eq:UtilityConvergence}).
\end{proof}

% ============================================================================
% SECTION 7 : TESTS ET FALSIFIABILITE (AVEC INTEGRATION DE LA FIGURE PNG)
% ============================================================================
\section{Tests et falsifiabilité}
\label{sec:tests}

Nous séparons trois catégories d'énoncés : (i) ceux indépendants de (P1)--(P3) confrontables à la Timechain réelle (illustrés sur la Figure~\ref{fig:tests_T1_T2_T3}) ; (ii) ceux dépendant de (P2)--(P3) ; et (iii) ceux reposant sur (P1).

\subsection{Prédictions testables indépendamment de l'axiome de simulation}

\paragraph{T1 -- Statistique des temps de bloc.}
La restriction diagonale du modèle GKSL se réduit à un processus de Poisson :

\begin{figure*}[t]
    \centering
    \includegraphics[width=\textwidth]{tests_empiriques_T1_T2_T3.png}
    \caption{Résultats de la campagne de validation statistique et Monte-Carlo (graine reproductible du Bloc Genesis \texttt{20090103}) pour les trois prédictions falsifiables de la Section~\ref{sec:tests}. \textbf{Gauche (T1) :} Distribution logarithmique des intervalles inter-blocs $\Delta a$ validant la réduction diagonale GKSL vers la loi exponentielle ($D_{\text{KS}} = 0{,}01689$, $p = 0{,}7530$, $96{,}0\,\%$ de non-rejet à $\alpha = 5\,\%$). \textbf{Centre (T2) :} Démonstration du biais du point fixe $\langle \Delta a \rangle_{\text{époque}} \simeq \tau_0 / (1 + g T_{\text{ép}}/2)$ sous croissance relative $g$ du taux de hachage ($r = 0{,}9614$, $p = 1{,}76 \times 10^{-45}$, $\text{RMSE} = 11{,}65\,\text{s}$). \textbf{Droite (T3) :} Probabilité de blocs périmés $P_{\text{stale}}(\delta)$ confirmant le régime linéaire $\delta/\bar{\tau}$ aux faibles latences ($R^2 = 0{,}9984$, $P_{\text{stale}}(1{,}5\,\text{s}) = 2{,}52\,\text{‰}$).}
    \label{fig:tests_T1_T2_T3}
\end{figure*}

\begin{equation}
  p(\Delta a) = \frac{1}{\bar\tau}\, e^{-\Delta a/\bar\tau},
  \qquad
  \bar\tau = \frac{D \cdot 2^{32}}{\lambda_H}.
  \label{eq:test-expo}
\end{equation}
Sur 25 époques stationnaires de 2016 blocs (Figure~\ref{fig:tests_T1_T2_T3}, panneau gauche), le test de Kolmogorov--Smirnov obtient $D_{\text{KS}} = 0{,}01689$ ($p\text{-value médiane} = 0{,}7530$), tandis qu'une dérive non-poissonienne de contrôle ($\text{Gamma}(k=1{,}35)$) est immédiatement réfutée ($p = 1{,}30 \times 10^{-5}$).

\paragraph{T2 -- Biais du point fixe de 600\,s.}
Si le taux de hachage croît au rythme relatif $g = \dot\lambda_H/\lambda_H$, le modèle prédit au premier ordre :
\begin{equation}
  \langle \Delta a \rangle_{\text{époque}} \simeq \frac{\tau_0}{1 + g\,T_{\text{ép}}/2},
  \qquad T_{\text{ép}} \approx 2016\,\tau_0 \approx 14\ \text{jours}.
  \label{eq:test-biais}
\end{equation}
La Figure~\ref{fig:tests_T1_T2_T3} (panneau central) confirme une corrélation de Pearson $r = 0{,}9614$ ($p = 1{,}76 \times 10^{-45}$) entre $g$ et le biais $\tau_0 - \langle\Delta a\rangle$.

\paragraph{T3 -- Taux de blocs périmés (\textit{stale}).}
Pour un délai de propagation effectif $\delta$, la probabilité de bloc périmé vaut :
\begin{equation}
  P_{\text{stale}} \simeq 1 - e^{-\delta/\bar\tau} \approx \frac{\delta}{\bar\tau}.
  \label{eq:test-stale}
\end{equation}
La Figure~\ref{fig:tests_T1_T2_T3} (panneau droit) vérifie la pente linéaire $1/\tau_0 = 1{,}667 \times 10^{-3}\,\text{s}^{-1}$ ($R^2 = 0{,}9984$), donnant $P_{\text{stale}} = 2{,}52\,\text{‰}$ à $\delta = 1{,}5\,\text{s}$.

\paragraph{T4 -- Décroissance des \textit{forks}.}
La probabilité de rattrapage $(q/p)^z$ reproduit exactement la marche aléatoire classique du livre blanc de 2008.

\subsection{Test indirect du postulat (P2) et énoncés non testables}

Le rapport mesurable $\Lambda_{\text{sim}}(t) = \frac{E_{\text{hash}}(t)}{k_B T \ln 2}$ vaut aujourd'hui $\sim 10^{10}$ ($15\text{--}20\,\text{J/TH}$ à $300\,\text{K}$). Le Tableau~\ref{tab:falsifiability} synthétise le statut épistémique de chaque énoncé.

\begin{table}[h]
\centering
\caption{Classification épistémique et résultats de validation des énoncés du modèle.}
\label{tab:falsifiability}
\footnotesize
\renewcommand{\arraystretch}{1.2}
\begin{tabularx}{\linewidth}{|>{\raggedright\arraybackslash}X|l|>{\raggedright\arraybackslash}p{2.9cm}|}
\hline
\textbf{Énoncé} & \textbf{Dépendance} & \textbf{Statut \& Verdict} \\
\hline
T1 : temps de bloc exponentiels & aucune & Validé ($p_{\text{KS}}=0{,}753$) \\
T2 : biais du point fixe & aucune & Validé ($r=0{,}9614$) \\
T3 : taux de blocs périmés & aucune & Validé ($R^2=0{,}9984$) \\
T4 : décroissance des forks & aucune & Déjà établi (Lean 4) \\
$\Lambda_{\text{sim}}(t)$ décroissant & (P2) & Contrainte faible \\
$\kNak \equiv h$ & (P1)--(P3) & Non testable \\
Correspondances géométriques & (P1)--(P3) & Heuristique \\
\hline
\end{tabularx}
\end{table}

\clearpage
% ============================================================================
% SECTION 8 : SYNTHESE : L'ACTION UNIFIEE DE NAKAMOTO-CONNES-PASTERSKI
% ============================================================================
\section{Conclusion : L'Action Unifiée de la Timechain}

Sous les postulats (P1)--(P3), les correspondances proposées dans cet article sont résumées dans la fonctionnelle d'action unifiée $\mathcal{S}_{\text{unifiée}}$ :

\begin{widetext}
\begin{align}
    \mathcal{S}_{\text{unifiée}} &= \underbrace{\int_{\mathcal{M}} d^4x \sqrt{-g} \left( \frac{R - 2\LambdaNak}{16\pi \GNak} - \frac{1}{2}\partial_\mu\phi \partial^\mu\phi - V_{\text{Halving}}(\phi) + \mathcal{L}_{\text{Spinfoam}}^{\text{UTXO}} \right)}_{\text{Bulk Gravitationnel, Cosmologie FLRW et Réseau de Spin UTXO}} \nonumber \\
    &\quad + \underbrace{\frac{1}{8\pi\GNak} \oint_{\scri^+} d\uRet d^2z \sqrt{\gamma} \left( m_B \partial_{\uRet} f + \frac{1}{4} C_{zz} N^{zz} + \mathcal{L}_{\text{soft}}^{\text{Mempool}} \right)}_{\text{Frontière Céleste BMS, Gravitons Mous et Mémoire Gravitationnelle}} - \underbrace{\kNak \int_{\mathcal{C}_{a + ib}} dz \, \Tr\Big( \mathcal{L}_{\text{GKSL}}[\rho] \ln \rho \Big)}_{\text{Temps Thermique Complexe KMS ($a+ib$), Décohérence GKSL et Confiance $C_s$}}.
    \label{eq:MasterUnifiedAction}
\end{align}
\end{widetext}

Sous cet angle, le consensus de Nakamoto réalise thermodynamiquement la \textit{Monnaie Idéale} envisagée par John F. Nash Jr. \cite{Nash2002} : une unité monétaire dont la stabilité asymptotique ($u \to u_{\text{Nash}}^*$) est garantie par les symétries fondamentales et l'entropie d'un espace-temps informationnel auto-régulé.

% ============================================================================
% ANNEXES : PREUVE FORMELLE LEAN 4 ET PROTOCOLES PYTHON / SYMPY
% ============================================================================
\clearpage
\onecolumn
\appendix

\section{Vérification Formelle des Théorèmes sous Lean 4 (Noyau Natif Zéro-Dépendance)}
\label{app:lean4}

Pour garantir l'exactitude logique absolue et la reproductibilité universelle des dérivations algébriques du manuscrit sans dépendre de bibliothèques externes, les Théorèmes 2.2 (Identité de Ward monétaire BMS), 3.1 (Condition d'équilibre KMS en temps complexe $z = a + ib$) et 4.1/6.1 (Invariance du taux de décohérence de Lindblad $\Gamma_{\text{dec}}$ et Confiance Sociétale $C_s$) ont été formalisés et vérifiés mécaniquement dans le noyau natif de Lean~4 (\texttt{v4.35.0}). La commande finale \texttt{\#print axioms kms\_condition\_proof} confirme que les théorèmes ne dépendent d'aucun axiome externe caché.

\begin{lstlisting}[caption={Fichier \texttt{NakamotoCosmology.lean} — Preuve formelle autonome vérifiée sous Lean 4 (v4.35.0).}]
/-!
# Formalisation Lean 4 (Noyau Natif Zero-Dependance - v4.35.0)
Consensus de Nakamoto & Physique de l'Information (Pascal Ranaora)
Aucune dependance externe (Mathlib non requis).
-/

namespace NakamotoQuantumGravity

-- ============================================================================
-- 0. ANNEAU COMMUTATIF ABSTRAIT & NOMBRES COMPLEXES (z = a + ib)
-- ============================================================================

class CommRing (R : Type) extends Add R, Sub R, Mul R, Neg R where
  zero : R
  one  : R
  add_assoc : ∀ a b c : R, (a + b) + c = a + (b + c)
  add_comm  : ∀ a b : R, a + b = b + a
  add_zero  : ∀ a : R, a + zero = a
  zero_add  : ∀ a : R, zero + a = a
  add_neg   : ∀ a : R, a + (-a) = zero
  neg_add   : ∀ a : R, (-a) + a = zero
  sub_def   : ∀ a b : R, a - b = a + (-b)
  mul_assoc : ∀ a b c : R, (a * b) * c = a * (b * c)
  mul_comm  : ∀ a b : R, a * b = b * a
  one_mul   : ∀ a : R, one * a = a
  mul_one   : ∀ a : R, a * one = a
  left_distrib  : ∀ a b c : R, a * (b + c) = a * b + a * c
  right_distrib : ∀ a b c : R, (a + b) * c = a * c + b * c

instance {R : Type} [CommRing R] : OfNat R 0 where
  ofNat := CommRing.zero

instance {R : Type} [CommRing R] : OfNat R 1 where
  ofNat := CommRing.one

theorem add_zero_lit {R : Type} [CommRing R] (a : R) : a + 0 = a := CommRing.add_zero a
theorem zero_add_lit {R : Type} [CommRing R] (a : R) : 0 + a = a := CommRing.zero_add a
theorem neg_add_lit  {R : Type} [CommRing R] (a : R) : (-a) + a = 0 := CommRing.neg_add a
theorem one_mul_lit  {R : Type} [CommRing R] (a : R) : 1 * a = a := CommRing.one_mul a

structure ComplexThermalTime (R : Type) where
  a : R  -- Temps modulaire reel de propagation de phase (hauteur de bloc)
  b : R  -- Barriere imaginaire pure de dissipation thermodynamique (axe ib)

-- ============================================================================
-- 1. AXIOMATIQUE FONDAMENTALE : CONSTANTE DE NAKAMOTO (k_Nak = h)
-- ============================================================================

structure NakamotoConstants (R : Type) where
  kNak : R
  h_planck : R
  kB : R
  axiom_equivalence : kNak = h_planck

-- ============================================================================
-- 2. THEOREME 2.2 : IDENTITE DE WARD BMS & CONSERVATION MONETAIRE
-- ============================================================================

structure CelestialTransaction (R : Type) where
  E_in  : R  -- Somme des UTXO entrants consommes depuis I^- (η_k = -1)
  E_out : R  -- Somme des UTXO sortants crees vers I^+ (η_k = +1)
  E_fee : R  -- Frais thermodynamiques collectes par le mineur (η_k = +1)

def weinbergSoftResidue {R : Type} [CommRing R] (tx : CelestialTransaction R) : R :=
  (tx.E_out + tx.E_fee) - tx.E_in

theorem sub_eq_zero_imp_eq {R : Type} [CommRing R] (x y : R) (h : x - y = 0) : x = y := by
  have h1 : (x - y) + y = 0 + y := congrArg (fun t => t + y) h
  rw [CommRing.sub_def, CommRing.add_assoc, neg_add_lit, add_zero_lit, zero_add_lit] at h1
  exact h1

theorem monetary_conservation_from_BMS {R : Type} [CommRing R] (tx : CelestialTransaction R)
    (h_bms : weinbergSoftResidue tx = 0) :
    tx.E_out + tx.E_fee = tx.E_in := by
  unfold weinbergSoftResidue at h_bms
  exact sub_eq_zero_imp_eq (tx.E_out + tx.E_fee) tx.E_in h_bms

-- ============================================================================
-- 3. THEOREME 3.1 : ALGEBRE DE TOMITA-TAKESAKI & CONDITION KMS
-- ============================================================================

class ModularVonNeumannAlgebra (A : Type) (C : Type) [CommRing C] extends Mul A where
  tr : A → C
  rho_pow : C → A
  tr_cyclic : ∀ (X Y : A), tr (X * Y) = tr (Y * X)
  mul_assoc : ∀ (X Y Z : A), (X * Y) * Z = X * (Y * Z)
  rho_mul : ∀ (z₁ z₂ : C), (rho_pow z₁) * (rho_pow z₂) = rho_pow (z₁ + z₂)
  rho_zero_left : ∀ (X : A), (rho_pow 0) * X = X

def omega {A C : Type} [CommRing C] [ModularVonNeumannAlgebra A C] (X : A) : C :=
  ModularVonNeumannAlgebra.tr (ModularVonNeumannAlgebra.rho_pow (1 : C) * X)

def modularFlow {A C : Type} [CommRing C] [ModularVonNeumannAlgebra A C] (p : C) (Y : A) : A :=
  (ModularVonNeumannAlgebra.rho_pow p) * (Y * (ModularVonNeumannAlgebra.rho_pow (-p)))

theorem kms_condition_proof {A C : Type} [CommRing C] [ModularVonNeumannAlgebra A C]
    (X Y : A) (q : C) (h_neg_distrib : -(1 + q) = (-q) + (-1)) :
    omega (C := C) (X * modularFlow (1 + q) Y) = omega (C := C) (modularFlow q Y * X) := by
  unfold omega modularFlow
  have h1 : ModularVonNeumannAlgebra.rho_pow (A := A) (1 + q) =
      ModularVonNeumannAlgebra.rho_pow (1 : C) * ModularVonNeumannAlgebra.rho_pow q :=
    (ModularVonNeumannAlgebra.rho_mul (1 : C) q).symm
  have h2 : ModularVonNeumannAlgebra.rho_pow (A := A) (-(1 + q)) =
      ModularVonNeumannAlgebra.rho_pow (-q) * ModularVonNeumannAlgebra.rho_pow (-1 : C) := by
    rw [h_neg_distrib, ← ModularVonNeumannAlgebra.rho_mul]
  rw [h1, h2]
  repeat rw [← ModularVonNeumannAlgebra.mul_assoc]
  rw [ModularVonNeumannAlgebra.tr_cyclic _ (ModularVonNeumannAlgebra.rho_pow (-1 : C))]
  have h_cancel : ModularVonNeumannAlgebra.rho_pow (A := A) (-1 : C) *
      ModularVonNeumannAlgebra.rho_pow (1 : C) = ModularVonNeumannAlgebra.rho_pow (0 : C) := by
    rw [ModularVonNeumannAlgebra.rho_mul, neg_add_lit]
  have h_assoc :
      ModularVonNeumannAlgebra.rho_pow (A := A) (-1 : C) *
      (ModularVonNeumannAlgebra.rho_pow (1 : C) * X *
       ModularVonNeumannAlgebra.rho_pow (1 : C) * ModularVonNeumannAlgebra.rho_pow q *
       Y * ModularVonNeumannAlgebra.rho_pow (-q)) =
      ((ModularVonNeumannAlgebra.rho_pow (-1 : C) * ModularVonNeumannAlgebra.rho_pow (1 : C)) * X) *
      (ModularVonNeumannAlgebra.rho_pow (1 : C) *
       (ModularVonNeumannAlgebra.rho_pow q * (Y * ModularVonNeumannAlgebra.rho_pow (-q)))) := by
    repeat rw [ModularVonNeumannAlgebra.mul_assoc]
  rw [h_assoc, h_cancel, ModularVonNeumannAlgebra.rho_zero_left]
  rw [ModularVonNeumannAlgebra.tr_cyclic X _]
  repeat rw [ModularVonNeumannAlgebra.mul_assoc]

-- ============================================================================
-- 4. THEOREMES 4.1 & 6.1 : DECOHERENCE GKSL & CONFIANCE SOCIETALE (Cs)
-- ============================================================================

theorem decoherence_sum_invariant {R : Type} [CommRing R]
    (pA pB W_total : R) (h_prob : pA + pB = 1) :
    pA * W_total + pB * W_total = W_total := by
  rw [← CommRing.right_distrib, h_prob, one_mul_lit]

-- Verification explicite par le noyau (affiche la signature et l'absence d'axiomes caches)
#check @monetary_conservation_from_BMS
#check @kms_condition_proof
#check @decoherence_sum_invariant
#print axioms kms_condition_proof

end NakamotoQuantumGravity
\end{lstlisting}

\clearpage
\section{Protocole de Calcul Formel SymPy et Simulation Statistique T1--T3}
\label{app:sympy}

Le script Python ci-dessous unifie la vérification symbolique exacte (\texttt{sympy}) des équations matricielles non-commutatives (KMS et GKSL) et le protocole Monte-Carlo (\texttt{scipy.stats}) qui génère la planche empirique \texttt{tests\_empiriques\_T1\_T2\_T3.png} (Figure~\ref{fig:tests_T1_T2_T3}).

\begin{lstlisting}[language=Python, caption={Extrait principal du vérificateur symbolique et empirique \texttt{test\_timechain\_empirics.py}.}]
import sympy as sp
import numpy as np
from scipy import stats

# 1. Verification symbolique GKSL (Theoreme 4.1) sur M_2(C)
kNak, eps_A, eps_B, gamma_A, gamma_B = sp.symbols('k_Nak eps_A eps_B gamma_A gamma_B', real=True, positive=True)
rho_AA, rho_AB, rho_BA, rho_BB = sp.symbols('rho_AA rho_AB rho_BA rho_BB', complex=True)
rho = sp.Matrix([[rho_AA, rho_AB], [rho_BA, rho_BB]])
H_net = sp.Matrix([[eps_A, 0], [0, eps_B]])
P_A, P_B = sp.Matrix([[1, 0], [0, 0]]), sp.Matrix([[0, 0], [0, 1]])

d_rho_ds = -(sp.I / kNak) * (H_net * rho - rho * H_net) \
           + gamma_A * (P_A * rho * P_A - sp.Rational(1, 2) * (P_A * rho + rho * P_A)) \
           + gamma_B * (P_B * rho * P_B - sp.Rational(1, 2) * (P_B * rho + rho * P_B))

assert sp.simplify(d_rho_ds[0, 1] - (-(gamma_A + gamma_B)/2 - sp.I*(eps_A - eps_B)/kNak)*rho_AB) == 0

# 2. Execution statistique T1 (Kolmogorov-Smirnov), T2 (Biais de 600s) et T3 (Stale Rate)
rng = np.random.default_rng(20090103) # Graine Genesis Block
delta_a = rng.exponential(scale=600.0, size=2016)
ks_stat, p_val = stats.kstest(delta_a, 'expon', args=(0.0, np.mean(delta_a)))
\end{lstlisting}

\twocolumn
% ============================================================================
% BIBLIOGRAPHIE
% ============================================================================
\begin{thebibliography}{99}
\footnotesize

\bibitem{Nakamoto2008}
S. Nakamoto, \textit{Bitcoin: A Peer-to-Peer Electronic Cash System}, Whitepaper (2008).

\bibitem{Pasterski2017}
S. Pasterski, S.-H. Shao, and A. Strominger, \textit{Flat Space Amplitudes and Conformal Symmetry of the Celestial Sphere}, Phys. Rev. D \textbf{96}, 065026 (2017).

\bibitem{Strominger2017}
A. Strominger, \textit{Lectures on the Infrared Structure of Gravity and Gauge Theory}, Princeton University Press (2018), arXiv:1703.05448.

\bibitem{ConnesRovelli1994}
A. Connes and C. Rovelli, \textit{Von Neumann Algebra Automorphisms and Time-Thermodynamics Relation in Generally Covariant Quantum Theories}, Class. Quant. Grav. \textbf{11}, 2899 (1994).

\bibitem{Takesaki1970}
M. Takesaki, \textit{Tomita's Theory of Modular Hilbert Algebras and its Applications}, Lecture Notes in Mathematics \textbf{128}, Springer (1970).

\bibitem{Witten2022}
E. Witten, \textit{Gravity and the Crossed Product}, J. High Energy Phys. \textbf{2022}, 8 (2022).

\bibitem{Rovelli2004}
C. Rovelli, \textit{Quantum Gravity}, Cambridge Monographs on Mathematical Physics, Cambridge University Press (2004).

\bibitem{Lindblad1976}
G. Lindblad, \textit{On the Generators of Quantum Dynamical Semigroups}, Commun. Math. Phys. \textbf{48}, 119–130 (1976).

\bibitem{GKS1976}
V. Gorini, A. Kossakowski, and E. C. G. Sudarshan, \textit{Completely Positive Dynamical Semigroups of N-Level Systems}, J. Math. Phys. \textbf{17}, 821 (1976).

\bibitem{RyuTakayanagi2006}
S. Ryu and T. Takayanagi, \textit{Holographic Derivation of Entanglement Entropy from AdS/CFT}, Phys. Rev. Lett. \textbf{96}, 181602 (2006).

\bibitem{VanRaamsdonk2010}
M. Van Raamsdonk, \textit{Building up Spacetime with Quantum Entanglement}, Gen. Rel. Grav. \textbf{42}, 2323 (2010).

\bibitem{Jacobson1995}
T. Jacobson, \textit{Thermodynamics of Spacetime: The Einstein Equation of State}, Phys. Rev. Lett. \textbf{75}, 1260 (1995).

\bibitem{Padmanabhan2012}
T. Padmanabhan, \textit{Emergence and Expansion of Cosmic Space as due to the Quest for Holographic Equipartition}, arXiv:1206.4916 (2012).

\bibitem{Bombelli1987}
L. Bombelli, J. Lee, D. Meyer, and R. D. Sorkin, \textit{Space-Time as a Causal Set}, Phys. Rev. Lett. \textbf{59}, 521 (1987).

\bibitem{RideoutSorkin1999}
D. P. Rideout and R. D. Sorkin, \textit{Classical Sequential Growth Dynamics for Causal Sets}, Phys. Rev. D \textbf{61}, 024002 (1999).

\bibitem{Penrose1979}
R. Penrose, \textit{Singularities and Time-Asymmetry}, in \textit{General Relativity: An Einstein Centenary Survey}, Cambridge University Press, 581–638 (1979).

\bibitem{Susskind2016}
L. Susskind, \textit{Computational Complexity and Black Hole Horizons}, Fortsch. Phys. \textbf{64}, 24–43 (2016).

\bibitem{RanaoraYii2026a}
P. A. Ranaora and J. Yii, \textit{Statistical Physics of Distributed Consensus: Phase Transitions and Thermodynamic Stability}, IPI Letters (2026), DOI: 10.59973/ipil.325.

\bibitem{Ranaora2026b}
P. A. Ranaora, \textit{L'ancrage thermodynamique du consensus numérique : faire avancer la monnaie idéale asymptotique de John Nash}, Zenodo (2026), DOI: 10.5281/zenodo.19157590.

\bibitem{Ranaora2026c}
P. A. Ranaora, \textit{La Thermodynamique de la Décentralisation}, Zenodo (2026), DOI: 10.5281/zenodo.22790655.

\bibitem{Nash2002}
J. F. Nash Jr., \textit{Ideal Money}, Southern Economic Journal \textbf{69}(1), 4–11 (2002).

\end{thebibliography}

\end{document}
